\section{What is a Pseudorandom Number Generator?}

Pseudorandom number generators (PRNGs) play a fundamental role in computer science, numerical methods, simulations, cryptography, randomized algorithms, and statistical experiments~\cite{knuth1997art,lecuyer2007testu01}. Their purpose is to generate sequences of numbers which, although produced by a deterministic algorithm, imitate certain properties of genuinely random sequences. In contrast to true randomness, which is usually associated with physical phenomena and inherently unpredictable processes, pseudorandomness is purely algorithmic. A PRNG is completely determined by its initial configuration. Therefore, once the initial value, called the seed, is fixed, the entire generated sequence can be reproduced exactly. This reproducibility is one of the main reasons why PRNGs are widely used in simulations, debugging, benchmarking, and experimental research. The central idea behind pseudorandom number generation is not to produce randomness in a philosophical or physical sense, but rather to produce deterministic sequences whose statistical behavior is sufficiently close to that of random sequences for a given application.

\subsection{Formal Definition}

Formally, a PRNG can be described as a discrete dynamical system consisting of:

\begin{itemize}
    \item a finite state space $S$,
    \item a deterministic transition function $f : S \to S$,
    \item an output function $g : S \to \mathbb{R}$.
\end{itemize}

The generator is defined by the recurrence relation:

$$
x_{n+1} = f(x_n),
$$

with initial state (seed) $x_0 \in S$, and output sequence:

$$
u_n = g(x_n).
$$

Here, $(x_n)$ represents the internal state sequence, while $(u_n)$ represents the generated pseudorandom numbers.

\subsection{Seed (Initial Value)}

The seed $x_0$ is the initial element of the state space $S$. It uniquely determines the entire sequence generated by the PRNG.

If the same seed is used multiple times, the generator produces the same sequence, which is a key property used for reproducibility in simulations, debugging, and numerical experiments.

Different seeds lead to different trajectories in the state space, increasing variability of outputs.

\subsection{State Space and Periodicity}

Since the state space $S$ is finite and the generator is deterministic, the sequence $(x_n)$ must eventually repeat. This follows from the pigeonhole principle: after a finite number of steps, some state must reoccur, causing the sequence to enter a cycle.

Thus, every PRNG becomes eventually periodic.

The period of a PRNG is the smallest positive integer $T$ such that:

$$
x_{n+T} = x_n \quad \text{for all } n \ge n_0,
$$

for some transient length $n_0 \ge 0$.

A good PRNG is characterized by a very long period, making repetition practically unobservable in applications.

\subsection{Uniformity and Equidistribution}

Uniformity refers to how evenly the output values are distributed over a given interval.

For normalized outputs $u_n \in [0,1]$, a sequence is considered uniformly distributed if, for any subinterval $[a,b] \subset [0,1]$, we have:

$$
\lim_{N \to \infty}
\frac{1}{N} \#\{1 \le n \le N : u_n \in [a,b]\}
= b - a.
$$

This property is also known as \textit{equidistribution}.

Lack of uniformity introduces bias and can significantly distort results in simulations and statistical applications, such as Monte Carlo methods.

\subsection{Computational Randomness}

Although PRNGs are deterministic, their goal is to produce sequences that are statistically indistinguishable from truly random sequences under practical tests.

In other words, while they are not random in a physical sense, they exhibit \textit{computational randomness}, meaning that without knowledge of the seed, the sequence appears random for algorithmic and statistical purposes.

\subsection{Motivation for Implementation}

A pseudorandom number generator can therefore be understood as a deterministic algorithm defined on a finite state space, designed to produce sequences with desirable statistical properties. A good general-purpose PRNG should satisfy several important requirements:

\begin{itemize}
    \item deterministic behavior, allowing reproducibility,
    \item a sufficiently long period,
    \item efficient computation,
    \item good uniformity and equidistribution properties,
    \item absence of easily detectable patterns,
    \item suitability for the intended application.
\end{itemize}

Many classical pseudorandom number generators can be described within this general framework. Examples include linear congruential generators, Lehmer generators, xorshift generators, and LFSR-based generators. They differ mainly in the choice of the transition function $f$, the structure of the state space $S$, and the output function $g$.

In the following sections, selected pseudorandom number generators will be implemented and analyzed in C++. Classical surveys and testing frameworks emphasize exactly this combination of theory, implementation and empirical validation~\cite{hull1962random,lecuyer2007testu01}. Before introducing particular algorithms, we first define a common interface that all generators will satisfy. This interface provides a uniform way to initialize generators, produce values, collect statistics, and generate visualizations.

\subsection{General Generator Interface}

In order to compare different pseudorandom number generators in a consistent way, all implementations in this project follow a common interface. This interface abstracts away the concrete internal mechanism of a generator and focuses only on the operations required to use and analyze it.

Conceptually, each generator is treated as an object with an internal state. The state can be initialized using a seed, advanced to produce the next output value, and described by basic information about its output range.

In the C++ implementation, this idea is expressed using a generic \texttt{RandomGenerator} concept. A type satisfying this concept must provide the following operations:

\begin{itemize}
    \item \texttt{next()}, which advances the generator and returns the next pseudorandom value,
    \item \texttt{seed(seed)}, which initializes or resets the internal state,
    \item \texttt{name()}, which returns the name of the generator,
    \item \texttt{min()} and \texttt{max()}, which describe the range of possible output values,
    \item \texttt{outputBits()}, which specifies how many bits of output are produced by a single call to \texttt{next()}.
\end{itemize}

This interface is directly connected to the mathematical model introduced earlier. The method \texttt{seed()} selects the initial state $x_0$, while \texttt{next()} corresponds to applying the transition function $f$ and then producing an observable output using the output function $g$.

Therefore, although different generators may use different recurrence relations and internal representations, they can still be tested and analyzed using the same code. This makes it possible to compare their periods, distributions, bit-level behavior, histograms, and correlations under a unified experimental framework.

Such an approach is also useful from a software engineering perspective. It separates the implementation of a particular generator from the code responsible for statistical analysis and visualization. As a result, adding a new PRNG requires only implementing the required interface, while the existing analysis tools can be reused without modification.

\cppfile[
    label={lst:random-generator-interface}
]{codes/RandomConcept.cpp}

